%%
\documentclass[%
]{ittmm}

\usepackage{polyglossia}
\setdefaultlanguage{russian}

%%% One can fix some overfulls
\sloppy

%% Minted listings support
%% Need pygment <http://pygments.org/> <http://pypi.python.org/pypi/Pygments>
\usepackage{minted}
%% auto break lines
\setminted{breaklines=true}

\addbibresource{main.bib}

%% end of the preamble, start of the body of the document source.
\begin{document}

%%
%% Rights management information.
%% CC-BY is default license.
\copyrightyear{2025}
\copyrightclause{Copyright for this paper by its authors.
  Use permitted under Creative Commons License Attribution 4.0
  International (CC BY 4.0).}

%%
%% This command is for the conference information
\conference{Information and Telecommunication Technologies and Mathematical Modeling of High-Tech Systems 2026 (ITTMM 2026), Moscow, April 13--18, 2026}

%%
%% The "title" command
\title{Градиентный бустинг для вероятностного прогнозирования LoS-блокировок в промышленных сценариях сетей 5G NR}

%% \tnotemark[1]
%% \tnotetext[1]{You can use this document as the template for preparing your publication. We recommend using the latest version of the ittmm style.}

%%
%% The "author" command and its associated commands are used to define
%% the authors and their affiliations.
\author{Д. А. Аветисян}[%
%% orcid=0000-0000-0000-0000,
email=yaedavid@ya.ru,
%% url=https://yaedavid.github.io,
]
\cormark[1]

\author{В. О. Бегишев}[%
%% orcid=0000-0000-0000-0000,
email=begishev-vo@rudn.ru,
]
%% \address[1]{RUDN University,
%%   6 Miklukho-Maklaya St, Moscow, 117198, Russian Federation}
%% \address[2]{Joint Institute for Nuclear Research,
%%   6 Joliot-Curie St, Dubna, 141980, Russian Federation}

%%% Footnotes
\cortext[1]{Автор, отвечающий за публикацию.}

%%
%% The abstract is a short summary of the work to be presented in the
%% article.
\begin{abstract}
    В работе рассматривается задача вероятностной оценки блокировки прямой видимости (LoS) в сетях 5G NR промышленного интернета вещей на основе методов машинного обучения. 
    Предложена модель градиентного бустинга над решающими деревьями для прогнозирования вероятности возникновения LoS-блокировки с учётом геометрических, радиофизических и средовых параметров. 
    Проведён численный анализ, показывающий эффективность предложенного подхода при моделировании сложной промышленной среды.
\end{abstract}

%%
%% Keywords. The author(s) should pick words that accurately describe
%% the work being presented. Separate the keywords with commas.
\begin{keywords}
  5G NR \sep
  промышленный интернет вещей \sep
  LoS-блокировка \sep
  машинное обучение \sep
  Gradient Boosting \sep
  вероятностное моделирование \sep
  радиоканал \sep
  mmWave
\end{keywords}

%% This command processes the author and affiliation and title
%% information and builds the first part of the formatted document.
\maketitle

\section{Введение}

Сети 5G NR миллиметрового диапазона (mmWave) рассматриваются как ключевая технологическая основа промышленного интернета вещей (IIoT). 
Высокая направленность и широкая полоса частот позволяют обеспечить требования eMBB и URLLC-приложений в автоматизированных производственных системах. 
Однако распространение сигнала в промышленной среде существенно осложняется явлением блокировки прямой видимости (Line-of-Sight, LoS), обусловленным наличием металлоконструкций, станков и подвижных механизмов.

В работе~\cite{Kondratyeva2022} предложена аналитическая модель вероятности LoS-блокировки для промышленных mmWave-развертываний 5G NR. 
Модель основана на сочетании фотограмметрического анализа прозрачности отдельных машин и стохастической геометрии (Manhattan Poisson Line Process), что позволяет получить выражения для вероятности блокировки в зависимости от плотности оборудования, геометрии размещения и высоты антенн.

Несмотря на аналитическую строгость, вычисление вероятности блокировки требует интегрирования по пространственным координатам и распределению высот машин, что усложняет применение модели в задачах оптимизации и системного планирования. 
В настоящей работе предлагается использовать методы машинного обучения для построения модели, аппроксимирующей аналитический вероятностный функционал.

\section{Аналитическая модель LoS-блокировки}

Рассматривается промышленное помещение, в котором производственные линии формируют стохастическую решётчатую структуру. Размещение линий моделируется Manhattan Poisson Line Process с интенсивностями $\lambda_x$ и $\lambda_y$ по осям.

Вероятность блокировки LoS для устройства, расположенного в точке $(l_x, l_y)$ определяется выражением вида

\begin{equation}
p_B(l_x,l_y)=\exp\left(-\kappa \nu (\Lambda_x l_x + \Lambda_y l_y)\right),
\end{equation}

где $\nu$ - вероятность занятости ячейки машиной, $\kappa$ - прозрачность машины, $\Lambda_x, \Lambda_y$ - эффективные интенсивности пересечений LoS с потенциальными блокирующими элементами.

В случае случайной высоты машин $H_B$ интенсивности задаются интегралами:

\begin{equation}
\Lambda_x=\int_0^{l_x} 
\lambda_x 
\left(
1-F_{H_B}\big(f_{L,x}(x)\big)
\right) 
dx,
\end{equation}

где $f_{L,x}(x)$ - высота луча LoS: 

\begin{equation}
f_{L,x}(x)=h_A-\frac{(h_A-h_U)x}{l_x}.
\end{equation}

Аналогичное выражение записывается для оси $y$.

Таким образом, итоговая вероятность блокировки представляет собой сложный нелинейный функционал от параметров среды:

\begin{equation}
p_B = f(d, h_A, h_U, \nu, \kappa, \lambda_x, \lambda_y, F_{H_B}).
\end{equation}

При многократных вычислениях (например, при оптимизации расположения базовых станций) непосредственное использование интегральной формы оказывается вычислительно затратным.

\section{Постановка задачи машинного обучения}

Целью является построение модели машинного обучения

\begin{equation}
\hat{f}(X) \approx p_B
\end{equation}

где вектор признаков

\begin{equation}
X = (d, h_A, h_U, \nu, \kappa, \lambda_x, \lambda_y, H_B)
\end{equation}

Задача формулируется как задача регрессии с ограниченным диапазоном выхода:

\begin{equation}
\hat{f}: \mathbb{R}^n \rightarrow [0,1]
\end{equation}

ML-модель должна воспроизводить поведение аналитической модели в широком диапазоне параметров.

% \subsection{Title variants}

% \mintinline{latex}{\title}
% command have the below options:
% \begin{itemize}
% \item
% \mintinline{latex}{title}
% : Document title. This is default option.
% \begin{minted}{latex}
% \title[mode=title]{This is a title}
% \end{minted}

% You can just omit it, like as follows:
% \begin{minted}{latex}
% \title{This is a title}
% \end{minted}

% \item
% \mintinline{latex}{alt}
% : Alternate title.
% \begin{minted}{latex}
% \title[mode=alt]{This is a alternate title}
% \end{minted}

% \item
% \mintinline{latex}{sub}
% : Sub title.
% \begin{minted}{latex}
% \title[mode=sub]{This is a sub title}
% \end{minted}
% You can just use
% \mintinline{latex}{\subtitle}
%  command, as follows:
% \begin{minted}{latex}
% \subtitle{This is a sub title}
% \end{minted}

% \item
% \mintinline{latex}{trans}
% : Translated title.
% \begin{minted}{latex}
% \title[mode=trans]{This is a translated title}
% \end{minted}

% \item
% \mintinline{latex}{transsub}
% : Translated sub title.
% \begin{minted}{latex}
% \title[mode=transsub]{This is a translated sub title}
% \end{minted}
% \end{itemize}

% \mintinline{latex}{\author}
% command have the below options:
% \begin{itemize}
% \item
% \mintinline{latex}{style}
%  : Style of author name (chinese)
% \item
% \mintinline{latex}{prefix}
%  : Prefix
% \item
% \mintinline{latex}{suffix}
%  : Suffix
% \item
% \mintinline{latex}{degree}
%  : Degree
% \item
% \mintinline{latex}{role}
%  : Role
% \item
% \mintinline{latex}{orcid}
%  : ORCID
% \item
% \mintinline{latex}{email}
%  : E-mail
% \item
% \mintinline{latex}{url}
%  : URL
% \end{itemize}

% Author names can have some kinds of marks and notes:
% \begin{itemize}
% \item affiliation mark:
% \mintinline{latex}{\author[<num>]}
% .
% \end{itemize}

% The author names and affiliations could be formatted in two ways:
% \begin{enumerate}
% \item Group the authors per affiliation.
% \item Use an explicit mark to indicate the affiliations.
% \end{enumerate}

% Author block example:
% \begin{minted}{latex}
% \author[1,2]{Author Name}[%
%     prefix=Prof.,
%     degree=D.Sc.,
%     role=Researcher,
%     orcid=0000-0000-000-0000,
%     email=name@example.com,
%     url=https://name.example.com
% ]

% \address[1]{Affiliation #1}
% \address[2]{Affiliation #2}
% \end{minted}

\section{Метод градиентного бустинга}

В качестве аппроксиматора используется ансамблевая модель градиентного бустинга:

\begin{equation}
F_M(x)=\sum_{m=1}^{M}\gamma_m h_m(x)
\end{equation}

где $h_m(x)$ — решающие деревья.

Минимизируется квадратичная функция потерь:

\begin{equation}
L=\sum_{i=1}^{N}\left(y_i-F(x_i)\right)^2
\end{equation}

Использован алгоритм Histogram-based Gradient Boosting, обеспечивающий высокую вычислительную эффективность при больших объёмах выборки.

Градиентный бустинг рассматривается как модель аналитического функционала, полученного в работе~\cite{Kondratyeva2022}.

\section{Численные эксперименты}

\subsection{Генерация датасета}

Датасет сформирован путём вычисления вероятности LoS-блокировки на основе аналитической модели из работы~\cite{Kondratyeva2022} при различных комбинациях параметров:

\begin{itemize}
  \item расстояние $d \in [5,100]$ м;
  \item плотность занятости $\nu \in [0.1,\,0.9]$;
  \item случайная высота машин;
  \item число наблюдений — до $200\,000$.
\end{itemize}

\subsection{Результаты in-domain}

Получены следующие метрики:

\begin{itemize}
  \item MAE $\approx 0.02$,
  \item RMSE $\approx 0.03$,
  \item $R^2 \approx 0.97$.
\end{itemize}

Модель точно аппроксимирует аналитический функционал.

\subsection{Out-of-distribution анализ}

Обучение проводилось при $\nu \le 0.6$, тестирование — при $\nu > 0.6$.

Результаты:

\begin{itemize}
  \item MAE $\approx 0.05$,
  \item RMSE $\approx 0.07$,
  \item $R^2 \approx 0.88$.
\end{itemize}

Даже при смене режима плотности модель демонстрирует устойчивость.

\begin{figure}[H]
\centering
\begin{subfigure}{0.48\textwidth}
    \centering
    \includegraphics[width=\linewidth]{figures/in_domain.png}
    \caption{In-domain тестирование}
\end{subfigure}
\hfill
\begin{subfigure}{0.48\textwidth}
    \centering
    \includegraphics[width=\linewidth]{figures/ood_test.png}
    \caption{OOD тестирование ($\nu > 0.6$)}
\end{subfigure}
\caption{Сравнение предсказанных и аналитических значений вероятности LoS-блокировки. Красная линия соответствует идеальному случаю $y=x$.}
\label{fig:predictions}
\end{figure}

\section{Обсуждение}

Предложенный подход:

1. Существенно ускоряет вычисление вероятности блокировки по сравнению с интегральной формой.  
2. Сохраняет физическую интерпретируемость.  
3. Устойчив к изменению плотности оборудования.  
4. Может использоваться в задачах оптимизации параметров сети.  

Таким образом, ML-подход расширяет аналитическую модель из работы~\cite{Kondratyeva2022}, обеспечивая вычислительно эффективную аппроксимацию.

\section{Заключение}

В работе предложена модель градиентного бустинга для вероятностного прогнозирования LoS-блокировок в промышленных сценариях 5G NR. 
Модель основана на аппроксимации аналитического стохастико-геометрического функционала, что обеспечивает сочетание физической корректности и вычислительной эффективности.

В дальнейшем планируется:

- учёт динамики подвижных механизмов,  
- расширение модели на временные ряды,  
- использование методов интерпретации (SHAP),  
- применение к реальным измерениям.  

%%
%% Define the bibliography file to be used
\printbibliography

%%
%% If your work has an appendix, this is the place to put it.
\appendix

\end{document}

%%% Local Variables:
%%% mode: LaTeX
%%% TeX-master: t
%%% End: